
    %%%%%%%%%%%%%%%%%%%%%%%%%%%%%%%%%%%%%%%%%%%%%%%%%%%%%%%%%%%%%%%%%%%%%%%%%%%%%%%%
    %%% Classe do documento
    \documentclass[12pt, addpoints]{exam}
    \UseRawInputEncoding

    %%%%%%%%%%%%%%%%%%%%%%%%%%%%%%%%%%%%%%%%%%%%%%%%%%%%%%%%%%%%%%%%%%%%%%%%%%%%%%%%
    % preambulo.tex
    %
    % Arquivo de configuração de packages para uso o LaTeX, conforme minhas
    % preferências e modelos pessoais.
    %
    % Para maiores informações, visite:
    %    
    %
    % NÃO ALTERE SE NÃO SOUBER O QUE ESTÁ FAZENDO!
    %%%%%%%%%%%%%%%%%%%%%%%%%%%%%%%%%%%%%%%%%%%%%%%%%%%%%%%%%%%%%%%%%%%%%%%%%%%%%%%%


    %%%%%%%%%%%%%%%%%%%%%%%%%%%%%%%%%%%%%%%%%%%%%%%%%%%%%%%%%%%%%%%%%%%%%%%%%%%%%%%%
    %%% Estruturas de controle:
    \input{static/utils/controles}

    %%%%%%%%%%%%%%%%%%%%%%%%%%%%%%%%%%%%%%%%%%%%%%%%%%%%%%%%%%%%%%%%%%%%%%%%%%%%%%%%
    %%% Compilação condicional em PDF:
    \input{static/utils/pdf}

    %%%%%%%%%%%%%%%%%%%%%%%%%%%%%%%%%%%%%%%%%%%%%%%%%%%%%%%%%%%%%%%%%%%%%%%%%%%%%%%%
    %%% Configurações de layout da página:
    \input{static/utils/layout}

    %%%%%%%%%%%%%%%%%%%%%%%%%%%%%%%%%%%%%%%%%%%%%%%%%%%%%%%%%%%%%%%%%%%%%%%%%%%%%%%%
    %%% Configurações para autores:
    \input{static/utils/autores}

    %%%%%%%%%%%%%%%%%%%%%%%%%%%%%%%%%%%%%%%%%%%%%%%%%%%%%%%%%%%%%%%%%%%%%%%%%%%%%%%%
    %%% Cabeçalho e rodapé:
    %%% Atenção! É MUITO DIFÍCIL ajustar apropriadamente os running headers
    %%% e running footers da primeira vez. Você pode confiar no LaTeX e deixar que
    %%% ele faça o ajuste automaticamente ou pode quebrar a cabeça e
    %%% tentar melhorar algo que o LaTeX já faz muito bem, mesmo que não fique
    %%% exatamente do jeito que você quer. O que você prefere? Se quiser ajustar
    %%% manualmente, descomente a linha abaixo e faça as configurações no arquivo
    %%% "utils/headings.tex".
    %\input{static/utils/headings}

    %%%%%%%%%%%%%%%%%%%%%%%%%%%%%%%%%%%%%%%%%%%%%%%%%%%%%%%%%%%%%%%%%%%%%%%%%%%%%%%%
    %%% Configuração para as fontes do documento:
    %%% Se você quiser usar fontes próprias ou do sistema, precisará ajustar as
    %%% configurações no arquivo "utils/fontes.tex".
    %%% ATENÇÃO: por padrão eu utilizo XeLaTeX com fontes proprietárias projetadas
    %%% por Matthew Butterick, adquiridas em https://mbtype.com, que não podem ser
    %%% distribuídas devido à licença de utilização. Se você usar XeLaTeX, você
    %%% DEVERÁ OBRIGATORIAMENTE ajustar as fontes no arquivo "utils/fontes.tex".
    \input{static/utils/fontes}

    %%%%%%%%%%%%%%%%%%%%%%%%%%%%%%%%%%%%%%%%%%%%%%%%%%%%%%%%%%%%%%%%%%%%%%%%%%%%%%%%
    %%% Sumário:
    \input{static/utils/toc}

    %%%%%%%%%%%%%%%%%%%%%%%%%%%%%%%%%%%%%%%%%%%%%%%%%%%%%%%%%%%%%%%%%%%%%%%%%%%%%%%%
    %%% Matemática:
    \input{static/utils/matematica}

    %%%%%%%%%%%%%%%%%%%%%%%%%%%%%%%%%%%%%%%%%%%%%%%%%%%%%%%%%%%%%%%%%%%%%%%%%%%%%%%%
    %%% Notas de rodapé:
    \input{static/utils/footnotes}

    %%%%%%%%%%%%%%%%%%%%%%%%%%%%%%%%%%%%%%%%%%%%%%%%%%%%%%%%%%%%%%%%%%%%%%%%%%%%%%%%
    %%% Referências cruzadas, links, citações:
    \input{static/utils/crossrefs}

    %%%%%%%%%%%%%%%%%%%%%%%%%%%%%%%%%%%%%%%%%%%%%%%%%%%%%%%%%%%%%%%%%%%%%%%%%%%%%%%%
    %%% Configurações para os hiperlinks em PDF:
    \input{static/utils/pdfconfig}

    %%%%%%%%%%%%%%%%%%%%%%%%%%%%%%%%%%%%%%%%%%%%%%%%%%%%%%%%%%%%%%%%%%%%%%%%%%%%%%%%
    %%% Glossário e índice remissivo:
    \input{static/utils/glossindex}

    %%%%%%%%%%%%%%%%%%%%%%%%%%%%%%%%%%%%%%%%%%%%%%%%%%%%%%%%%%%%%%%%%%%%%%%%%%%%%%%%
    %%% Referências bibliográficas:
    \input{static/utils/bibliografia}

    %%%%%%%%%%%%%%%%%%%%%%%%%%%%%%%%%%%%%%%%%%%%%%%%%%%%%%%%%%%%%%%%%%%%%%%%%%%%%%%%
    %%% Cores:
    \input{static/utils/cores}

    %%%%%%%%%%%%%%%%%%%%%%%%%%%%%%%%%%%%%%%%%%%%%%%%%%%%%%%%%%%%%%%%%%%%%%%%%%%%%%%%
    %%% Computação:
    \input{static/utils/computacao}

    %%%%%%%%%%%%%%%%%%%%%%%%%%%%%%%%%%%%%%%%%%%%%%%%%%%%%%%%%%%%%%%%%%%%%%%%%%%%%%%%
    %%% Gráficos:
    \input{static/utils/graficos}

    %%%%%%%%%%%%%%%%%%%%%%%%%%%%%%%%%%%%%%%%%%%%%%%%%%%%%%%%%%%%%%%%%%%%%%%%%%%%%%%%
    %%% Ícones:
    \input{static/utils/icones}

    %%%%%%%%%%%%%%%%%%%%%%%%%%%%%%%%%%%%%%%%%%%%%%%%%%%%%%%%%%%%%%%%%%%%%%%%%%%%%%%%
    %%% Blocos:
    \input{static/utils/blocos}

    %%%%%%%%%%%%%%%%%%%%%%%%%%%%%%%%%%%%%%%%%%%%%%%%%%%%%%%%%%%%%%%%%%%%%%%%%%%%%%%%
    %%% Boxes coloridos:
    \input{static/utils/colorbox}

    %%%%%%%%%%%%%%%%%%%%%%%%%%%%%%%%%%%%%%%%%%%%%%%%%%%%%%%%%%%%%%%%%%%%%%%%%%%%%%%%
    %%% Tabelas:
    \input{static/utils/tabelas}

    %%%%%%%%%%%%%%%%%%%%%%%%%%%%%%%%%%%%%%%%%%%%%%%%%%%%%%%%%%%%%%%%%%%%%%%%%%%%%%%%
    %%% Ambientes de listas:
    \input{static/utils/listas}

    %%%%%%%%%%%%%%%%%%%%%%%%%%%%%%%%%%%%%%%%%%%%%%%%%%%%%%%%%%%%%%%%%%%%%%%%%%%%%%%%
    %%% Ambientes floats e similares:
    \input{static/utils/floats}

    %%%%%%%%%%%%%%%%%%%%%%%%%%%%%%%%%%%%%%%%%%%%%%%%%%%%%%%%%%%%%%%%%%%%%%%%%%%%%%%%
    %%% Ajustes para a classe EXAM:
    \input{static/utils/exam}

    %%%%%%%%%%%%%%%%%%%%%%%%%%%%%%%%%%%%%%%%%%%%%%%%%%%%%%%%%%%%%%%%%%%%%%%%%%%%%%%%
    %%% Ajustes para textos:
    \input{static/utils/texto}

    %%%%%%%%%%%%%%%%%%%%%%%%%%%%%%%%%%%%%%%%%%%%%%%%%%%%%%%%%%%%%%%%%%%%%%%%%%%%%%%%
    %%% Ambientes, aliases, comandos e customizações em geral para serem
    %%% utilizadas nos documentos. Defina aqui qualquer customização específica
    %%% para seus documentos!
    \input{static/utils/customizacoes}

    %%%%%%%%%%%%%%%%%%%%%%%%%%%%%%%%%%%%%%%%%%%%%%%%%%%%%%%%%%%%%%%%%%%%%%%%%%%%%%%%
    %%% Ajuste do layout, espaçamento de linhas e etc.:
    \geometry{a4paper, portrait, left=2.0cm, right=2.0cm, top=2.0cm, bottom=2.0cm}
    \singlespacing

    %%%%%%%%%%%%%%%%%%%%%%%%%%%%%%%%%%%%%%%%%%%%%%%%%%%%%%%%%%%%%%%%%%%%%%%%%%%%%%%%
    %%% Configurações de cabeçalho e rodapé (não use fancyhdr pois ocorre conflito):
    \pagestyle{headandfoot}
    \runningheadrule
    \firstpageheader{}{}{}
    \runningheader{Sem disciplina}{}{Avaliação}
    \firstpagefooter{}{}{Página \thepage\ de \numpages}
    \runningfooter{}{}{Página \thepage\ de \numpages}
    \runningfootrule

    %%%%%%%%%%%%%%%%%%%%%%%%%%%%%%%%%%%%%%%%%%%%%%%%%%%%%%%%%%%%%%%%%%%%%%%%%%%%%%%%
    %%% Configurações para as propriedades do PDF:
    \hypersetup{
    hidelinks,           % Comente para web, descomente para impressão
    colorlinks=true,      % True para web, False para impressão
    pdftitle=Minha nova Prof: Avaliação,
    pdfauthor=Professor,
    pdfsubject=Sem disciplina,
    pdfkeywords=['Sem tag'],
    pdfinfo={
        CreationDate={}, % Ex.: D:AAAAMMDDHH24MISS
        ModDate={}       % Ex.: D:AAAAMMDDHH24MISS
    }
    }
    \input{static/utils/pdfconfig}

    %%%%%%%%%%%%%%%%%%%%%%%%%%%%%%%%%%%%%%%%%%%%%%%%%%%%%%%%%%%%%%%%%%%%%%%%%%%%%%%%
    %%% Começa o documento
    \begin{document}

    %%%%%%%%%%%%%%%%%%%%%%%%%%%%%%%%%%%%%%%%%%%%%%%%%%%%%%%%%%%%%%%%%%%%%%%%%%%%%%%%
    %%% Folha de instruções padronizadas da UVV para a prova
    \pdfbookmark[1]{Instruções}{instrucoes}
    \begin{figure}[H]
        \begin{center}
            \includegraphics[scale=0.3]{{static/imgs/logo-uvv.png}}
        \end{center}	
    \end{figure}

    \vspace{-1cm}

    \begin{table}[H]
        \centering
        \begin{tabular}{|llll|}
            \hline
            \multicolumn{2}{|l|}{\textbf{Disciplina:} Sem disciplina} & 
            \multicolumn{1}{l|}{\multirow{3}{*}{\textbf{\makecell{Nota:\\ \,\\ \,}}}} & 
            \multirow{3}{*}{\textbf{\makecell{Coordenador:\\ \,\\ \,}}} \\ 
            \cline{1-2}
            \multicolumn{2}{|l|}{\textbf{Professor:} Professor \hspace{4.72cm} } & 
            \multicolumn{1}{l|}{} & \\ 
            \cline{1-2}
            \multicolumn{2}{|l|}{\textbf{Aluno:}} & 
            \multicolumn{1}{l|}{} & \\ 
            \hline
            \multicolumn{1}{|l|}{\textbf{Turma:} \hspace{6.3cm} } & 
            \multicolumn{1}{l|}{\textbf{Semestre:}} & 
            \multicolumn{2}{l|}{\textbf{Valor:} 10 pontos} \\ 
            \hline
            \multicolumn{1}{|l|}{\textbf{Data:}} & 
            \multicolumn{3}{l|}{\textbf{Avaliação:}} \\ 
            \hline
        \end{tabular}
    \end{table}

    \begin{center}
        \fbox{\setlength\fboxsep{0.3cm} \fbox{\parbox{15.4cm}{
            \begin{center}
                \textbf{INSTRUÇÕES PARA A REALIZAÇÃO DESTA AVALIAÇÃO:}
            \end{center}
            \begin{itemize}[leftmargin=*]
                \item Esta é a primeira avaliação da disciplina Sem disciplina, 
                    e engloba o conteúdo referente Sem tag.
                \item Esta avaliação contém 22 questões objetivas, \textbf{todas obrigatórias}.
                \item \textbf{Leia atentamente cada questão!} As questões objetivas terão uma,
                    e apenas uma única, resposta a ser marcada.
                \item \textbf{Desligue o celular} antes de começar. A avaliação é
                    \textbf{individual} e \textbf{sem consulta}.
                \item As questões podem ser respondidas, na prova, com lápis ou caneta. Ao final
                    da prova \textbf{{VOCÊ DEVERÁ PREENCHER O GABARITO COM CANETA
                    DE COR PRETA}} \textbf{\underline{OU AZUL}} para que seja feita a correção
                    automatizada através da leitura do padrão de respostas marcado no
                    gabarito.
                \item \textbf{CUIDADO ao preencher o gabarito} pois eles são \textbf{INDIVIDUAIS
                    e NOMEADOS} (cada estudante tem um gabarito próprio específico, com um
                    código de identificação único). \textbf{Questões rasuradas são anuladas}
                    pelo software de leitura do gabarito.
                \item Ao final da prova, \textbf{DEVOLVA A PROVA E O GABARITO} para o professor.
                \item Siga todas as normas de \textbf{Integridade Acadêmica} da disciplina.
                    Alunos que forem flagrados com qualquer espécie de ``cola'' ou trocando
                    informações com outros alunos terão suas avaliações recolhidas, as notas
                    zeradas, e a situação será encaminhada para a coordenação para a aplicação
                    das penalidades previstas pela universidade.
                \item Com 90 minutos de avaliação você terá tempo de até 4,min por questão,
                    em média.
                \item Boa avaliação!
            \end{itemize}
            \vspace{2cm}
        }}}
    \end{center}
    \newpage

    %%%%%%%%%%%%%%%%%%%%%%%%%%%%%%%%%%%%%%%%%%%%%%%%%%%%%%%%%%%%%%%%%%%%%%%%%%%%%%%%
    %%% Inicia o bloco de questões:
    \begin{questions}

    \newpage
    \question

Considere uma cena representada no sistema de referência do universo (SRU), uma \textit{window} definida pelo par de coordenadas (0,0)- (100,100) e uma \textit{viewport} definida pelo par de coordenadas (20,30)- (300,100). Considere ainda que as coordenadas que definem \textit{window} e \textit{viewport} correspondem, respectivamente, aos limites inferior esquerdo e superior direito de ambas. Analise as afirmativas abaixo levando em consideração os conceitos clássicos de \textit{window} e \textit{viewport} e assinale a alternativa correta.



\begin{itemize}

    \item[I)] \textit{Window} e \textit{viewport} estão definidas no SRU.

    \item[II)] No processo de mapeamento desta \textit{window} para esta \textit{viewport} haverá modificação na relação de aspecto.

    \item[III)] O mapeamento da \textit{window} redefinida pelo par de coordenadas (0,0) – (50,50) para a mesma \textit{viewport} (20,30) – (300,100) corresponde a uma operação de \textit{zoom out} sobre o mesmo universo.

\end{itemize}
\begin{choices}
\choice As alternativas I e III são falsas.
\choice As alternativas I e II são verdadeiras.
\choice As afirmativas II e III são verdadeiras.
\choice As alternativas I e II são falsas.
\choice Apenas a afirmativa III é verdadeira.
\end{choices}

\question

Redes semânticas, frames e lógica são formalismos utilizados principalmente em: 
\begin{choices}
\choice IA distribuída
\choice inferência em sistemas especialistas
\choice descoberta de conhecimento em bases de dados
\choice redes neurais
\choice representação de conhecimento
\end{choices}

\question

Dentre as características do modelo relacional e do modelo de objetos em bancos de dados, qual afirmação é \textbf{INCORRETA}?


\begin{choices}
\choice O modelo de objetos é mais adequado para a representação de tipos abstratos de dados.
\choice O relacionamento de herança é diretamente representado no modelo relacional.
\choice O modelo de objetos provê uma representação bem próxima de linguagens de programação.
\choice O modelo de objetos possui mais recursos estruturais para a representação de dados que o relacional.
\choice O relacionamento binário \( N \times M \) é representado de modo semelhante nos dois modelos.	
\end{choices}

\question

Qual das seguintes igualdades é verdadeira?



\begin{enumerate}

    \item[(I)] \( n^2 = \mathcal{O}(n^3) \)

    \item[(II)] \( 2 \cdot n + 1 = \mathcal{O}(n^2) \)

    \item[(III)] \( n^3 = \mathcal{O}(n^2) \)

    \item[(IV)] \( 3 \cdot n + 5 \cdot n \log n = \mathcal{O}(n) \)

    \item[(V)] \( \log n + \sqrt{n} = \mathcal{O}(n) \)

\end{enumerate}



As opções para a resposta correta são:
\begin{choices}
\choice somente III, IV e V
\choice somente I, III e IV
\choice somente I e II
\choice somente II, III e IV
\choice somente I, II e V
\end{choices}

\question

Considere o algoritmo da busca sequencial de um elemento em um conjunto com \( n \) elementos. A expressão que representa o tempo médio de execução desse algoritmo para uma busca bem sucedida é:
\begin{choices}
\choice \( n(n + 1)/2 \)
\choice \( (n + 1)/2 \)
\choice \( \log_2 n \)
\choice \( n^2 \)
\choice \( n \log n \)
\end{choices}

\question

O pipeline de visualização de objetos tridimensionais reúne um conjunto de transformações e processos aplicados a primitivas geométricas. Sobre essas transformações e processos pode-se dizer que:

\begin{enumerate}[label=\Roman*.]

    \item Os objetos devem corresponder a sólidos.

    \item As coordenadas dos vértices sofrem transformação de acordo com a posição e orientação do observador.

    \item Um volume de visualização correspondente a um paralelepípedo é determinado pela adoção de projeção perspectiva.

    \item A fase final do pipeline corresponde à rasterização dos polígonos.

\end{enumerate}

Selecione a alternativa correta:
\begin{choices}
\choice Apenas a afirmativa IV está verdadeira.
\choice Todas as afirmativas são verdadeiras.
\choice Apenas as afirmativas I e III são falsas.
\choice Apenas a afirmativa IV é falsa.
\choice As afirmativas II e III são falsas.
\end{choices}

\question

Uma definição lógica, abstrata e auto-contida dos objetos que modelam a

estrutura de armazenamento de dados, dos operadores que modelam o comportamento

e dos demais conceitos que modelam a integridade, e que, juntos, constituem a

máquina abstrata com a qual os usuários interagem é chamada de:


\begin{choices}
\choice Modelo de dados
\choice Modelo relacional
\choice Banco de dados
\choice Modelo entidade-relacionamento
\choice Nenhuma das respostas acima
\end{choices}

\question

Três atletas \( A \), \( B \) e \( C \) competiram, ao pares, numa corrida de \( d \) metros. Considerando que cada atleta teve o mesmo desempenho (ou seja, a mesma velocidade) ao competir com adversários distintos, e sabendo-se que



\begin{itemize}

    \item \( A \) venceu \( B \) chegando 20 metros à frente

    \item \( B \) venceu \( C \) chegando 10 metros à frente

    \item \( A \) venceu \( C \) chegando 28 metros à frente,

\end{itemize}



podemos afirmar que a corrida tem
\begin{choices}
\choice 150 metros
\choice 100 metros
\choice 200 metros
\choice 50 metros
\choice 110 metros
\end{choices}

\question

A interposição de um circuito de memória cache entre o processador e a memória principal (RAM)
\begin{choices}
\choice aumenta o tráfego de instruções e/ou dados entre memória e disco
\choice diminui o tráfego de instruções e/ou dados no barramento de memória
\choice diminui o tráfego de instruções e/ou dados entre memória e disco
\choice permite acessos concorrentes à memória RAM
\choice aumenta o tráfego de instruções e/ou dados no barramento de memória
\end{choices}

\question Assinale a alternativa \textbf{incorreta}:
\begin{choices}
\choice Modelo de dados é a mesma coisa que o diagrama relacional
\choice Um SGBD não tem interface gráfica
\choice O modelo relacional de dados foi criado por Edgar Frank Cood
\choice O MariaDB é um SGBD
\choice No modelo relacional, tudo o que o usuário percebe são tabelas
\end{choices}

\question

Seja a árvore binária abaixo a representação de um espaço de estados para um problema \( p \), em que o estado inicial é \( a \), e \( i \) e \( f \) são estados finais.



\begin{figure}[H]

    \centering

    \includegraphics[width=0.5\linewidth]{static/imgs/defaul.png}

    \caption{Questão 59}

    \label{fig:enter-label}

\end{figure}



Um algoritmo de busca em largura-primeiro forneceria a seguinte sequência de estados como primeira alternativa a um caminho-solução para o problema \( p \):
\begin{choices}
\choice \( a \ b \ e \ i \)
\choice \( a \ c \ f \)
\choice \( a \ b \ d \ h \ e \ i \)
\choice \( a \ b \ d \ e \ f \)
\choice \( a \ b \ c \ d \ e \ f \)
\end{choices}

\question

São motivos importantes e gerais para o estudo das linguagens de programação,

exceto:
\begin{choices}
\choice Aprender a sintaxe e a semântica de uma determinada linguagem
\choice Melhorar o embasamento na escolha das linguagens para um projeto
\choice Avanço geral da computação
\choice Aumentar a capacidade de expressar idéias
\choice Entender a importância da implementação
\end{choices}

\question

Qual das informações a seguir \textbf{NÃO} é colocada no registro de ativação na chamada de funções?


\begin{choices}
\choice Valor de retorno da função
\choice Link para a subrotina chamadora
\choice Endereço de retorno
\choice Variáveis locais estáticas
\choice Estado dos registradores
\end{choices}

\question

Quais dos algoritmos de ordenação abaixo possuem tempo no pior caso e tempo médio de execução proporcional a \( O(n \log n) \).
\begin{choices}
\choice Merge sort e heap sort
\choice Merge sort e bubble sort
\choice Quicksort e merge sort
\choice Heap sort e selection sort
\choice Bubble sort e Quick sort
\end{choices}

\question

Uma característica de uma arquitetura RISC é:
\begin{choices}
\choice Possui um pequeno conjunto de instruções simples
\choice Possui um grande conjunto de instruções complexas
\choice O processador é formado por válvulas e transistores
\choice Uma arquitetura de alto risco pois o mercado de hardware evolui muito rapidamente
\choice A arquitetura é constituída de milhares de processadores
\end{choices}

\question

A sentença descritiva resultante de um processo de mensuração, correspondendo a

um fato bruto observado/conhecido que pode ser registrado, que não foi

processado para revelar seu significado/importância, que tem um certo

significado implícito e que deve estar corretamente formatado para o

armazenamento, processamento e apresentação é:
\begin{choices}
\choice Conhecimento
\choice Dado
\choice Sabedoria
\choice Informação
\choice Processamento
\end{choices}

\question

Marque a alternativa onde todos os conceitos estão corretos.


\begin{choices}
\choice Em um diagrama de fluxo de dados, uma entidade externa representa um produtor ou um consumidor de informação e está fora dos limites do sistema modelado; cada processo pode ser refinado, para explicitar um maior detalhamento; um DFD contém dois níveis de detalhamento; um processo é um transformador de informação e também está fora do sistema; o nível 0 de um DFD representa o sistema como um todo e indica os principais usuários e as funções do sistema.
\choice Em um diagrama de fluxo de dados uma entidade externa representa uma fonte ou destino das informações processadas pelo sistema e está fora dos limites do sistema modelado; cada processo pode ser refinado, para explicitar um maior detalhamento; um DFD pode conter vários níveis de detalhamento; um processo é um transformador de informação e também está fora do sistema; o nível 0 de um DFD representa o sistema como um todo e indica as principais fontes e destinos das informações.
\choice Em um diagrama de fluxo de dados uma entidade externa representa uma fonte ou destino das informações processadas pelo sistema e está fora dos limites do sistema modelado; cada processo pode ser refinado, para explicitar um maior detalhamento; um DFD pode conter vários níveis de detalhamento; um processo é um transformador de informação; o nível 0 de um DFD representa o sistema como um todo e indica as principais fontes e destinos das informações, usualmente referenciado por Diagrama de Contexto.
\choice Nenhuma das alternativas anteriores.
\choice Em um diagrama de fluxo de dados uma entidade externa representa um produtor ou um consumidor de informação e está fora dos limites do sistema modelado; cada processo deve ser refinado, para explicitar um maior detalhamento; um DFD pode conter vários níveis de detalhamento; um processo é um transformador de informação e também está fora do sistema; o nível 0 de um DFD representa o sistema como um todo e indica os principais usuários e as funções do sistema.
\end{choices}

\question

Considere as seguintes afirmações:

\begin{enumerate}[label=\Roman*.]

    \item A classe de linguagens aceita por um Autômato Finito Determinístico (AFD) não é a mesma que um Autômato Finito Não Determinístico (AFND).

    \item Para algumas expressões regulares não é possível construir um AFD.

    \item A expressão regular \((b + ba)^+\) aceita as "strings" de b's e a's começando com b e não tendo dois a's consecutivos.

\end{enumerate}

Selecione a afirmativa correta:
\begin{choices}
\choice Apenas a afirmativa III é verdadeira.
\choice As afirmativas II e III são falsas.
\choice As afirmativas I e II são verdadeiras.
\choice As afirmativas I e III são verdadeiras.
\choice As afirmativas I e III são falsas.
\end{choices}

\question

Um banco de dados OLAP (Online Analytical Processing --- Processamento Analítico

Online) é um banco de dados organizado de forma a dar suporte ao Business

Intelligence (BI --- Inteligência de Negócio) que, de forma simplificada, é o

processo de analisar dados para obter informações que podem ser utilizadas na

tomada de decisões comerciais e estratégicas. Em relação aos bancos de dados

OLAP, marque a afirmação correta:
\begin{choices}
\choice São otimizados para o processamento de grande volume de dados históricos e métricas de negócio para a tomada de decisões estratégicas, permitindo e facilitando a análise de dados para a geração de informações e conhecimentos para a tomada de decisões.
\choice São otimizados para o processamento de grande volume de transações diárias, o que permite a análise e a tomada de decisões em tempo real, favorecendo assim a administração empresarial.
\choice São preparados para responder perguntas do tipo ``Quais alunos estão matriculados na turma de Bancos de Dados I neste semestre?'
\choice São otimizados para consulta e relatórios de rotina de dados.
\choice Geralmente os dados de origem do OLTP são de bancos de dados OLAP.
\end{choices}

\question

Por que não podemos utilizar apelidos de coluna na cláusula WHERE da SQL?
\begin{choices}
\choice Porque só podemos utilizar aliases, e não apelidos.
\choice Porque a cláusula SELECT é avaliada antes da WHERE.
\choice Nós podemos utilizar apelidos na cláusula WHERE.
\choice Porque a cláusula WHERE é avaliada antes da SELECT.
\choice Porque os apelidos só são avaliados depois da cláusula ORDER BY.
\end{choices}

\question

Na linguagem SQL, qual das seguintes palavras-chave é usada para modificar os

dados em uma tabela existente?


\begin{choices}
\choice INSERT
\choice CREATE
\choice SELECT
\choice UPDATE
\choice Nenhuma das alternativas acima
\end{choices}

\question

A função abaixo, escrita na linguagem C, quando executada para \( n = 5 \), faz quantas chamadas recursivas (excluindo a primeira chamada da função)? 

\begin{verbatim} int fat (int n) 

{ 

if (n == 1) return n; 

else return (n*fat(n-1)); 

} 

\end{verbatim}
\begin{choices}
\choice 5
\choice 6
\choice 0
\choice 1
\choice 4
\end{choices}



    %%%%%%%%%%%%%%%%%%%%%%%%%%%%%%%%%%%%%%%%%%%%%%%%%%%%%%%%%%%%%%%%%%%%%%%%%%%%%%%%
    %%% Finaliza o bloco de questões:
    \end{questions}
    \end{document}
    